\section*{摘要}

\noindent
动态多跳无线网络的时分多址（TDMA）调度需要同时兼顾吞吐量、公平性与拓扑扰动恢复能力，
而传统启发式调度依赖手工权值难以适应拓扑与业务的快速变化，
直接将申请概率替换为深度强化学习（DRL）策略又会破坏已有的安全退避机制并产生饥饿节点。
本文提出\textbf{密度感知服务债务 PPO}（Density-Aware Service-Debt PPO, DASD-PPO）调度框架，
其核心思路是将 PPO 策略以\emph{乘数}形式叠加在启发式申请概率之上，
并在采样阶段引入两跳安全掩膜、服务债务推力、请求预算约束与密度自适应控制，
使学习型策略在拓扑感知归纳偏置下做出公平且鲁棒的逐时隙决策。
DASD-PPO 采用 LSTM Actor--Critic 主干（每节点独立、约 13.8K 参数），
对每个数据时隙单独输出 Bernoulli 申请概率，并配合逐时隙奖励分解解决多维动作的信用归因问题；
仿真器与训练器通过 POSIX 命名管道完成异步在线交互。
在 OMNeT++ 6.3 平台上的对比实验涵盖两个静态主场景（行人 MANET 与 UAV 稀疏拓扑）
与四类动态扰动场景（边切换、桥接断裂、节点重接入、节点永久脱离），
将 DASD-PPO 与 PPO-Base、HC-PPO 及七种启发式/论文启发式近似算法做了同口径对比。
结果显示，DASD-PPO 在 UAV 稀疏拓扑下将饥饿比例从 \(0.5556\) 显著降至 \(\mathbf{0.2500}\)
（全表最低），在 Bridge Break 扰动下取得最小 PDR 衰减，
在 Node Rejoin 恢复阶段取得最高 throughput（\(2.6072\)），
学习型方法整体在 PDR 与 throughput 上数倍优于论文启发式近似基线，
证实在线 RL 调度对动态拓扑具有结构性鲁棒优势。

\vspace{0.5em}
\noindent\textbf{关键词：}时分多址；强化学习；MAC 协议；近端策略优化；空间复用；公平性；服务债务；在线学习
